%======================================================================
\section{Central}\label{sec:central}
%======================================================================
A central recebe o que os hubs publicam, guarda o histórico, mostra os quartos e avisa o
cuidador. Ela não decide nada: as regras e os alertas chegam prontos das placas. Tudo roda em
contêineres Docker, descritos num único arquivo do Docker Compose, num PC ou num Raspberry Pi 4
(Tabela~\ref{tab:central}).

\begin{table}[H]
\centering
\caption{Serviços da central.}
\label{tab:central}
\small
\begin{tabular}{@{}l>{\raggedright\arraybackslash}p{11.8cm}@{}}
\toprule
Serviço & Função \\
\midrule
Mosquitto & broker MQTT com usuário e senha por cliente e ACL por tópico (Seção~\ref{sec:mqtt}) \\
Telegraf & assina \code{idoso/\#} e grava cada mensagem no InfluxDB, com o site, o quarto e o identificador como \textit{tags}; sem código próprio \\
InfluxDB & banco de séries temporais, com retenção de 30 dias \\
Grafana & painel ``Quartos'', com seletor de quarto; o painel é versionado em JSON no repositório \\
Notificador & programa em Python que assina \code{idoso/+/+/alerta} e envia cada alerta ao Telegram do cuidador \\
\midrule
Simulador & programa em Python que faz o papel de um hub, para testar a central sem hardware \\
\code{demo\_pulseira.py} & demonstração da pulseira com uma placa só, pela USB (Seção~\ref{sec:demo}) \\
\bottomrule
\end{tabular}
\end{table}

\subsection{Painel}

O painel ``Quartos'' mostra, para o quarto escolhido: o estado do hub (online ou offline, pelo
\textit{Last Will}), a presença, a temperatura atual e o histórico, as portas, o estado da
pulseira, a linha do tempo de portas e presença, a tabela de alertas (com a probabilidade de
queda dada pela pulseira), a bateria e o nível de atividade da pulseira
(Figura~\ref{fig:grafana}). As consultas foram testadas com o simulador; a aparência com dados
reais ainda vai ser conferida.

\begin{figure}[H]
\centering
\fotoouespaco{grafana_quartos}{0.95\linewidth}{captura do painel ``Quartos'' do Grafana com dados do hub (ou do simulador)}
\caption{Painel ``Quartos'' no Grafana.}
\label{fig:grafana}
\end{figure}

\subsection{Notificador}

O notificador transforma cada alerta numa mensagem legível (tipo, quarto e horário) e a manda
pelo Telegram. Ele foi feito para não perder alertas:
\begin{itemize}
  \item a sessão MQTT é persistente: os alertas publicados enquanto ele estava fora do ar
        chegam quando ele volta;
  \item se o Telegram não aceitar a mensagem (sem internet, por exemplo), ele espera e tenta de
        novo, com espera crescente até 60\,s; só uma mensagem recusada pelo próprio Telegram
        é descartada;
  \item \code{id} e \code{estado} repetidos (QoS 1) são descartados;
  \item um alerta com dados inesperados vira uma mensagem simples, em vez de derrubar o
        programa.
\end{itemize}

\begin{figure}[H]
\centering
\fotoouespaco{telegram_alerta}{0.42\linewidth}{captura do celular com um alerta de queda no Telegram}
\caption{Alerta de queda chegando ao cuidador pelo Telegram.}
\label{fig:telegram}
\end{figure}

\subsection{Demonstração com uma placa}\label{sec:demo}

Para mostrar o modelo de queda funcionando sem montar o sistema todo, a pulseira tem um perfil
de demonstração: sem hub, com o modo de coleta ligado e a serial a 921\,600\,baud. No
computador, o \code{demo\_pulseira.py} lê a serial e mostra, ao vivo, o $\amag$ filtrado, cada
disparo do gatilho com a probabilidade dada pelo modelo e uma faixa ``QUEDA DETECTADA'' quando
a pulseira decide queda. Opcionalmente, ele faz o papel do hub e publica o alerta no broker,
no mesmo formato, e o alerta aparece no Grafana e no Telegram. Se o link USB cair, a leitura
reabre a porta sozinha.

\begin{figure}[H]
\centering
\fotoouespaco{demo_pulseira}{0.85\linewidth}{janela do demo\_pulseira.py durante uma queda simulada}
\caption{A janela da demonstração com uma placa.}
\label{fig:demo}
\end{figure}
